% =========================================================================
% Capítulo 0: Repaso
% Tema 02: Cálculo vectorial
% Problema 07: Hélice y triedro de Frenet-Serret
% Origen: Ayudantía 24, 2026
% =========================================================================
\subsubsection{Hélice y triedro de Frenet-Serret}
Una partícula de masa $m$ describe una trayectoria espiral ascendente (hélice circular recta). Su vector de posición respecto a un marco inercial fijo $\mathcal{N}$ está dado en coordenadas cartesianas por:
$$\vec{r}(t) = R \cos(\omega t)\vec{e}_x + R \sin(\omega t)\vec{e}_y + c t \vec{e}_z$$
donde $R$, $\omega$ y $c$ son constantes positivas. Determine:
\begin{enumerate}
    \item El triedro de Frenet-Serret $(\vec{e}_t, \vec{e}_n, \vec{e}_b)$, la velocidad y la aceleración.
    \item La ecuación del plano osculador.
    \item La curvatura $\kappa$, la torsión $\tau$ y la verificación de las fórmulas de Serret-Frenet.
\end{enumerate}

% --- SOLUCIÓN ---------------------------------------------------------------
\begin{solucion}{Hélice y triedro de Frenet-Serret}
\textbf{Marcos de referencia y Sistemas de coordenadas}
\begin{itemize}
  \item $\{\vec{e}_x,\vec{e}_y,\vec{e}_z\}$: Marco inercial cartesiano fijo $\mathcal{N}$.
  \item $\{\vec{e}_t,\vec{e}_n,\vec{e}_b\}$: Marco móvil intrínseco (Triedro de Frenet-Serret).
\end{itemize}

\begin{center}
    % Dibujo 3D de la hélice y el triedro
    \begin{tikzpicture}[>=Latex, x={(0.7cm, 0.4cm)}, y={(-0.8cm, 0.4cm)}, z={(0cm, 1.2cm)}, scale=1.3]
        % Ejes cartesianos
        \draw[->, gray!80, thick] (0,0,0) -- (3.5,0,0) node[right] {$x$};
        \draw[->, gray!80, thick] (0,0,0) -- (0,3.5,0) node[left] {$y$};
        \draw[->, gray!80, thick] (0,0,0) -- (0,0,3) node[above] {$z$};

        % Hélice 3D
        \draw[thick, blue, smooth, samples=100, domain=0:4.5*pi] 
            plot ({2*cos(\x r)}, {2*sin(\x r)}, {0.15*\x});

        % Punto P (evaluado en un t especifico, ej. 2.5 pi)
        \coordinate (P) at ({2*cos(2.5*pi r)}, {2*sin(2.5*pi r)}, {0.15*2.5*pi});
        \filldraw[black] (P) circle (1.5pt) node[right, xshift=2pt] {$m$};

        % Triedro de Frenet (vectores)
        % Tangente (derivada de la curva)
        \draw[->, thick, red] (P) -- ++({-1.5*sin(2.5*pi r)}, {1.5*cos(2.5*pi r)}, {0.5}) node[above right] {$\vec{e}_t$};
        % Normal (hacia el eje z)
        \draw[->, thick, green!60!black] (P) -- ++({-1.5*cos(2.5*pi r)}, {-1.5*sin(2.5*pi r)}, {0}) node[below left] {$\vec{e}_n$};
        % Binormal (producto cruz)
        \draw[->, thick, orange] (P) -- ++({0.5*cos(2.5*pi r)}, {0.5*sin(2.5*pi r)}, {1.2}) node[above left] {$\vec{e}_b$};
    \end{tikzpicture}
\end{center}

\subsubsection*{Coordenadas intrínsecas: cinemática y plano osculador}
\textbf{Velocidad y Vector Tangente ($\vec{e}_t$):}
Derivando el vector posición respecto al tiempo:
$$\vec{v}(t) = \dot{\vec{r}}(t) = -R\omega\sin(\omega t)\vec{e}_x + R\omega\cos(\omega t)\vec{e}_y + c\vec{e}_z$$
La rapidez (módulo de la velocidad) es constante:
$$v = \|\vec{v}\| = \sqrt{(-R\omega\sin(\omega t))^2 + (R\omega\cos(\omega t))^2 + c^2} = \sqrt{R^2\omega^2 + c^2} \equiv v_0$$
El vector unitario tangente se define como la dirección de la velocidad:
$$\vec{e}_t = \frac{\vec{v}}{v_0} = \frac{1}{v_0}\Big(-R\omega\sin(\omega t)\vec{e}_x + R\omega\cos(\omega t)\vec{e}_y + c\vec{e}_z \Big)$$

\textbf{Aceleración y Vector Normal ($\vec{e}_n$):}
Derivando la velocidad respecto al tiempo:
$$\vec{a}(t) = \dot{\vec{v}}(t) = -R\omega^2\cos(\omega t)\vec{e}_x - R\omega^2\sin(\omega t)\vec{e}_y + 0\vec{e}_z$$
En coordenadas intrínsecas, sabemos que $\vec{a} = \dot{v}\vec{e}_t + \frac{v^2}{\rho}\vec{e}_n$. Como $v_0$ es constante, la aceleración tangencial es nula ($\dot{v}=0$). Toda la aceleración es centrípeta (normal):
$$a = \|\vec{a}\| = R\omega^2 \quad \Rightarrow \quad \vec{e}_n = \frac{\vec{a}}{a} = -\cos(\omega t)\vec{e}_x - \sin(\omega t)\vec{e}_y$$

\textbf{Vector Binormal ($\vec{e}_b$):}
Completa el triedro ortonormal diestro mediante el producto cruz $\vec{e}_b = \vec{e}_t \times \vec{e}_n$:
$$\vec{e}_b = \frac{1}{v_0}\Big( c\sin(\omega t)\vec{e}_x - c\cos(\omega t)\vec{e}_y + R\omega\vec{e}_z \Big)$$

\textbf{Plano Osculador:}
El plano osculador contiene instantáneamente a los vectores $\vec{e}_t$ y $\vec{e}_n$, por lo que su vector normal es $\vec{e}_b$. Su ecuación para un punto genérico $(X,Y,Z)$ en un instante $t$ se define imponiendo que el vector relativo sea ortogonal a $\vec{e}_b$:
\begin{align*}
    (\vec{R} - \vec{r}(t)) \cdot \vec{e}_b &= 0  \\
    c\sin(\omega t)\big[X - R\cos(\omega t)\big] - c\cos(\omega t)\big[Y - R\sin(\omega t)\big] + R\omega\big[Z - ct\big] &= 0
\end{align*}


\subsubsection*{Curvatura, torsión y fórmulas de Serret-Frenet}
Las fórmulas de Serret-Frenet describen cómo rota el triedro respecto a la longitud de arco $s$. Usando la regla de la cadena, $\frac{d}{ds} = \frac{1}{v_0} \frac{d}{dt}$:
\begin{align*}
    \dot{\vec{e}}_t &= v_0 \kappa \vec{e}_n \tag{I} \\
    \dot{\vec{e}}_n &= -v_0 \kappa \vec{e}_t + v_0 \tau \vec{e}_b \tag{II} \\
    \dot{\vec{e}}_b &= -v_0 \tau \vec{e}_n \tag{III}
\end{align*}

\textbf{Cálculo de la Curvatura ($\kappa$):}
Usando la expresión intrínseca de la aceleración normal $a_n = \kappa v^2$:
$$R\omega^2 = \kappa (v_0)^2$$
\respuesta{$\displaystyle\kappa = \frac{R\omega^2}{R^2\omega^2 + c^2}$}
\textit{Verificación Serret-Frenet (I):} 
$\dot{\vec{e}}_t = \frac{1}{v_0}(-R\omega^2\cos(\omega t)\vec{e}_x - R\omega^2\sin(\omega t)\vec{e}_y) = \left(\frac{R\omega^2}{v_0}\right) \vec{e}_n$. 
Como $v_0 \kappa = \sqrt{R^2\omega^2 + c^2} \left(\frac{R\omega^2}{R^2\omega^2 + c^2}\right) = \frac{R\omega^2}{v_0}$, la ecuación se cumple.

\textbf{Cálculo de la Torsión ($\tau$):}
La torsión mide cuánto se sale la curva del plano osculador. La calculamos derivando $\vec{e}_b$ en el tiempo y usando la fórmula (III):
$$\dot{\vec{e}}_b = \frac{1}{v_0} \Big( c\omega\cos(\omega t)\vec{e}_x + c\omega\sin(\omega t)\vec{e}_y \Big) = -\left(\frac{c\omega}{v_0}\right) \vec{e}_n$$
Igualando con la ecuación (III) que indica $\dot{\vec{e}}_b = -v_0 \tau \vec{e}_n$:
$$-v_0 \tau = -\frac{c\omega}{v_0}$$
\respuesta{$\displaystyle\tau = \frac{c\omega}{v_0^2} = \frac{c\omega}{R^2\omega^2 + c^2}$}

\textit{Conclusión Geométrica:} Dado que $R$, $\omega$ y $c$ son constantes, la trayectoria tiene \textbf{curvatura constante} (la espiral no se ensancha) y \textbf{torsión constante} (el "paso" o elevación de la rosca es uniforme), siendo un ejemplo clásico del teorema de Lancret.
\end{solucion}
